% Options for packages loaded elsewhere
\PassOptionsToPackage{unicode}{hyperref}
\PassOptionsToPackage{hyphens}{url}
%
\documentclass[
  ignorenonframetext,
  aspectratio=169,
  russian,
]{beamer}
\usepackage{pgfpages}
\setbeamertemplate{caption}[numbered]
\setbeamertemplate{caption label separator}{: }
\setbeamercolor{caption name}{fg=normal text.fg}
\beamertemplatenavigationsymbolshorizontal
% Prevent slide breaks in the middle of a paragraph
\widowpenalties 1 10000
\raggedbottom
\setbeamertemplate{part page}{
  \centering
  \begin{beamercolorbox}[sep=16pt,center]{part title}
    \usebeamerfont{part title}\insertpart\par
  \end{beamercolorbox}
}
\setbeamertemplate{section page}{
  \centering
  \begin{beamercolorbox}[sep=12pt,center]{section title}
    \usebeamerfont{section title}\insertsection\par
  \end{beamercolorbox}
}
\setbeamertemplate{subsection page}{
  \centering
  \begin{beamercolorbox}[sep=8pt,center]{subsection title}
    \usebeamerfont{subsection title}\insertsubsection\par
  \end{beamercolorbox}
}
\AtBeginPart{
  \frame{\partpage}
}
\AtBeginSection{
  \ifbibliography
  \else
    \frame{\sectionpage}
  \fi
}
\AtBeginSubsection{
  \frame{\subsectionpage}
}

\usepackage{amsmath,amssymb}
\usepackage{iftex}
\ifPDFTeX
  \usepackage[T1]{fontenc}
  \usepackage[utf8]{inputenc}
  \usepackage{textcomp} % provide euro and other symbols
\else % if luatex or xetex
  \usepackage{unicode-math}
  \defaultfontfeatures{Scale=MatchLowercase}
  \defaultfontfeatures[\rmfamily]{Ligatures=TeX,Scale=1}
\fi
\usepackage{lmodern}
\usefonttheme{serif} % use mainfont rather than sansfont for slide text
\ifPDFTeX\else  
    % xetex/luatex font selection
    \setmainfont[Ligatures=TeX]{Liberation Serif}
    \setsansfont[Ligatures=TeX,Scale=MatchLowercase]{Liberation Sans}
    \setmonofont[Scale=MatchLowercase,Scale=0.9]{Liberation Mono}
\fi
% Use upquote if available, for straight quotes in verbatim environments
\IfFileExists{upquote.sty}{\usepackage{upquote}}{}
\IfFileExists{microtype.sty}{% use microtype if available
  \usepackage[]{microtype}
  \UseMicrotypeSet[protrusion]{basicmath} % disable protrusion for tt fonts
}{}
\usepackage{xcolor}
\newif\ifbibliography
\setlength{\emergencystretch}{3em} % prevent overfull lines
\setcounter{secnumdepth}{-\maxdimen} % remove section numbering


\providecommand{\tightlist}{%
  \setlength{\itemsep}{0pt}\setlength{\parskip}{0pt}}\usepackage{longtable,booktabs,array}
\usepackage{calc} % for calculating minipage widths
\usepackage{caption}
% Make caption package work with longtable
\makeatletter
\def\fnum@table{\tablename~\thetable}
\makeatother
\usepackage{graphicx}
\makeatletter
\newsavebox\pandoc@box
\newcommand*\pandocbounded[1]{% scales image to fit in text height/width
  \sbox\pandoc@box{#1}%
  \Gscale@div\@tempa{\textheight}{\dimexpr\ht\pandoc@box+\dp\pandoc@box\relax}%
  \Gscale@div\@tempb{\linewidth}{\wd\pandoc@box}%
  \ifdim\@tempb\p@<\@tempa\p@\let\@tempa\@tempb\fi% select the smaller of both
  \ifdim\@tempa\p@<\p@\scalebox{\@tempa}{\usebox\pandoc@box}%
  \else\usebox{\pandoc@box}%
  \fi%
}
% Set default figure placement to htbp
\def\fps@figure{htbp}
\makeatother
% definitions for citeproc citations
\NewDocumentCommand\citeproctext{}{}
\NewDocumentCommand\citeproc{mm}{%
  \begingroup\def\citeproctext{#2}\cite{#1}\endgroup}
\makeatletter
 % allow citations to break across lines
 \let\@cite@ofmt\@firstofone
 % avoid brackets around text for \cite:
 \def\@biblabel#1{}
 \def\@cite#1#2{{#1\if@tempswa , #2\fi}}
\makeatother
\newlength{\cslhangindent}
\setlength{\cslhangindent}{1.5em}
\newlength{\csllabelwidth}
\setlength{\csllabelwidth}{3em}
\newenvironment{CSLReferences}[2] % #1 hanging-indent, #2 entry-spacing
 {\begin{list}{}{%
  \setlength{\itemindent}{0pt}
  \setlength{\leftmargin}{0pt}
  \setlength{\parsep}{0pt}
  % turn on hanging indent if param 1 is 1
  \ifodd #1
   \setlength{\leftmargin}{\cslhangindent}
   \setlength{\itemindent}{-1\cslhangindent}
  \fi
  % set entry spacing
  \setlength{\itemsep}{#2\baselineskip}}}
 {\end{list}}
\usepackage{calc}
\newcommand{\CSLBlock}[1]{\hfill\break\parbox[t]{\linewidth}{\strut\ignorespaces#1\strut}}
\newcommand{\CSLLeftMargin}[1]{\parbox[t]{\csllabelwidth}{\strut#1\strut}}
\newcommand{\CSLRightInline}[1]{\parbox[t]{\linewidth - \csllabelwidth}{\strut#1\strut}}
\newcommand{\CSLIndent}[1]{\hspace{\cslhangindent}#1}

\IfFileExists{plex-otf.sty}{
  %% Full TeXlive
  % \usepackage[%
  %   % math,
  %   RM={Scale=0.94},SS={Scale=0.94},SScon={Scale=0.94},TT={Scale=MatchLowercase,FakeStretch=0.9},DefaultFeatures={Ligatures=Common}
  % ]{plex-otf}
}{
  %% TinyTeX
  \usepackage{libertine}
}

%%% Load theme
% https://deic.uab.cat/~iblanes/beamer_gallery/
\IfFileExists{beamerthemegotham.sty}{
  %% Full TeXlive
  \usetheme{gotham}
  \gothamset{
    numbering=totalpagenumber,
    parttocframe default=off,
    sectiontocframe default=off,
    subsectiontocframe default=off,
  }
}{
  %% TinyTeX
  \usetheme{Madrid}
}
\makeatletter
\@ifpackageloaded{caption}{}{\usepackage{caption}}
\AtBeginDocument{%
\ifdefined\contentsname
  \renewcommand*\contentsname{Содержание}
\else
  \newcommand\contentsname{Содержание}
\fi
\ifdefined\listfigurename
  \renewcommand*\listfigurename{Список иллюстраций}
\else
  \newcommand\listfigurename{Список иллюстраций}
\fi
\ifdefined\listtablename
  \renewcommand*\listtablename{Список таблиц}
\else
  \newcommand\listtablename{Список таблиц}
\fi
\ifdefined\figurename
  \renewcommand*\figurename{Рисунок}
\else
  \newcommand\figurename{Рисунок}
\fi
\ifdefined\tablename
  \renewcommand*\tablename{Таблица}
\else
  \newcommand\tablename{Таблица}
\fi
}
\@ifpackageloaded{float}{}{\usepackage{float}}
\floatstyle{ruled}
\@ifundefined{c@chapter}{\newfloat{codelisting}{h}{lop}}{\newfloat{codelisting}{h}{lop}[chapter]}
\floatname{codelisting}{Список}
\newcommand*\listoflistings{\listof{codelisting}{Листинги}}
\makeatother
\makeatletter
\makeatother
\makeatletter
\@ifpackageloaded{caption}{}{\usepackage{caption}}
\@ifpackageloaded{subcaption}{}{\usepackage{subcaption}}
\makeatother

\ifLuaTeX
\usepackage[bidi=basic]{babel}
\else
\usepackage[bidi=default]{babel}
\fi
\babelprovide[main,import]{russian}
\ifPDFTeX
\else
\babelfont{rm}[Ligatures=TeX]{Liberation Serif}
\fi
\babelprovide[import]{english}
% get rid of language-specific shorthands (see #6817):
\let\LanguageShortHands\languageshorthands
\def\languageshorthands#1{}
\usepackage{csquotes}
\usepackage{bookmark}

\IfFileExists{xurl.sty}{\usepackage{xurl}}{} % add URL line breaks if available
\urlstyle{same} % disable monospaced font for URLs
\hypersetup{
  pdftitle={Выполнение индивидуального проекта},
  pdfauthor={Пашутина Анна Алексеевна},
  pdflang={ru-RU},
  hidelinks,
  pdfcreator={LaTeX via pandoc}}


\title{Выполнение индивидуального проекта}
\subtitle{Этап 5}
\author{Пашутина Анна Алексеевна}
\date{2026-05-16}

\begin{document}
\frame{\titlepage}

\renewcommand*\contentsname{Содержание}
\begin{frame}[allowframebreaks]
  \frametitle{Содержание}
  \setcounter{tocdepth}{1}
  \tableofcontents
\end{frame}

\section{Информация}\label{ux438ux43dux444ux43eux440ux43cux430ux446ux438ux44f}

\begin{frame}{Докладчик}
\phantomsection\label{ux434ux43eux43aux43bux430ux434ux447ux438ux43a}
\begin{columns}[c]
\begin{column}{0.7\linewidth}
\begin{itemize}[<+->]
\tightlist
\item
  Пашутина Анна Алексеевна
\item
  Студентка
\item
  Российский университет дружбы народов
\item
  1032253642@rudn.ru
\item
  группа НПИ-Бл-02-25
\end{itemize}
\end{column}

\begin{column}{0.3\linewidth}
\end{column}
\end{columns}
\end{frame}

\begin{frame}{Цель}
\phantomsection\label{ux446ux435ux43bux44c}
\begin{itemize}[<+->]
\tightlist
\item
  Создать индивидуальный сайт, постепенно его заполняя.
\end{itemize}
\end{frame}

\begin{frame}{Задание}
\phantomsection\label{ux437ux430ux434ux430ux43dux438ux435}
\begin{itemize}[<+->]
\tightlist
\item
  Добавить с сайту все остальные элементы.
\item
  Сделать записи для персональных проектов.
\item
  Сделать пост по прошедшей неделе.
\item
  Добавить пост на тему по выбору: языки научного программирования.
\end{itemize}
\end{frame}

\begin{frame}{Оформление проекта}
\phantomsection\label{ux43eux444ux43eux440ux43cux43bux435ux43dux438ux435-ux43fux440ux43eux435ux43aux442ux430}
\begin{itemize}[<+->]
\tightlist
\item
  Добавим в файл index.md нашего проекта ссылку на свой гитхаб.
\end{itemize}

\includegraphics[width=0.7\linewidth,height=\textheight,keepaspectratio]{./image/image1.png}
\end{frame}

\begin{frame}{Фото для ссылки}
\phantomsection\label{ux444ux43eux442ux43e-ux434ux43bux44f-ux441ux441ux44bux43bux43aux438}
\begin{itemize}[<+->]
\tightlist
\item
  Создаю аватарку для ссылки на проект гитхаб на сайте.
\end{itemize}

\includegraphics[width=0.7\linewidth,height=\textheight,keepaspectratio]{./image/image2.png}
\end{frame}

\begin{frame}{Пост по прошедшей неделе}
\phantomsection\label{ux43fux43eux441ux442-ux43fux43e-ux43fux440ux43eux448ux435ux434ux448ux435ux439-ux43dux435ux434ux435ux43bux435}
\begin{itemize}[<+->]
\tightlist
\item
  Пишем пост о прошедшей неделе.
\end{itemize}

\includegraphics[width=0.7\linewidth,height=\textheight,keepaspectratio]{./image/image3.png}
\end{frame}

\begin{frame}{Пост о языках программирования}
\phantomsection\label{ux43fux43eux441ux442-ux43e-ux44fux437ux44bux43aux430ux445-ux43fux440ux43eux433ux440ux430ux43cux43cux438ux440ux43eux432ux430ux43dux438ux44f}
\begin{itemize}[<+->]
\tightlist
\item
  Пишем пост о языках программирования: какой лучше выбрать новичку?
\end{itemize}

\includegraphics[width=0.7\linewidth,height=\textheight,keepaspectratio]{./image/image4.png}
\end{frame}

\begin{frame}{Ссылка на гитхаб-проект на сайте}
\phantomsection\label{ux441ux441ux44bux43bux43aux430-ux43dux430-ux433ux438ux442ux445ux430ux431-ux43fux440ux43eux435ux43aux442-ux43dux430-ux441ux430ux439ux442ux435}
\begin{itemize}[<+->]
\tightlist
\item
  Так выглядит ссылка на мой гитхаб-проект на моем сайте.
\end{itemize}

\includegraphics[width=0.7\linewidth,height=\textheight,keepaspectratio]{./image/image5.png}
\end{frame}

\begin{frame}{Пост о языках прораммирования на сайте}
\phantomsection\label{ux43fux43eux441ux442-ux43e-ux44fux437ux44bux43aux430ux445-ux43fux440ux43eux440ux430ux43cux43cux438ux440ux43eux432ux430ux43dux438ux44f-ux43dux430-ux441ux430ux439ux442ux435}
\begin{itemize}[<+->]
\tightlist
\item
  Так выглядит пост о языках программирования на моем сайте.
\end{itemize}

\includegraphics[width=0.7\linewidth,height=\textheight,keepaspectratio]{./image/image6.png}
\end{frame}

\begin{frame}{Пост по прошедшей неделе на сайте}
\phantomsection\label{ux43fux43eux441ux442-ux43fux43e-ux43fux440ux43eux448ux435ux434ux448ux435ux439-ux43dux435ux434ux435ux43bux435-ux43dux430-ux441ux430ux439ux442ux435}
\begin{itemize}[<+->]
\tightlist
\item
  Так выглядит пост о прошедшей неделе на моем сайте.
\end{itemize}

\includegraphics[width=0.7\linewidth,height=\textheight,keepaspectratio]{./image/image7.png}
\end{frame}

\section{Выводы}\label{ux432ux44bux432ux43eux434ux44b}

\begin{frame}{Выводы}
\begin{itemize}[<+->]
\tightlist
\item
  В результате работы были добавлены проекты и новые посты.
\end{itemize}
\end{frame}

\section*{Список
литературы}\label{ux441ux43fux438ux441ux43eux43a-ux43bux438ux442ux435ux440ux430ux442ux443ux440ux44b}
\addcontentsline{toc}{section}{Список литературы}

\begin{frame}{Список литературы}
\phantomsection\label{refs}
\begin{CSLReferences}{0}{1}
\end{CSLReferences}
\end{frame}




\end{document}
